\setcounter{section}{19}

  \zwischenueberschrift{Übungsaufgaben}

\inputaufgabe
{}
{

Zeige, dass die
\definitionsverweis {zweite Fundamentalmatrix}{}{}
\definitionsverweis {symmetrisch}{}{}
ist.

}
{} {}

\inputaufgabe
{}
{

Es sei
\maabbeledisp {\psi} {\R^2} {\R
} {(x,y)} { x^2-xy-y^3
} {,}
und sei $Y$ der
\definitionsverweis {Graph}{}{}
zu $\psi$ mit dem nach oben gerichteten
\definitionsverweis {Einheitsnormalenfeld}{}{}
und
\maabbdisp {\varphi  =
\operatorname{Id}  \times \psi} {\R^2} {Y
} {}
die Parametrisierung des Graphen. Bestimme die
\definitionsverweis {zweite Fundamentalmatrix}{}{}
zu $\varphi$.

}
{} {}

\inputaufgabe
{}
{

Es sei
\maabbeledisp {\psi} {\R^2} {\R
} {(x,y)} { e^{xy} -x^2y^3  \sin  \left(  x-y^2 \right)
} {,}
und sei $Y$ der
\definitionsverweis {Graph}{}{}
zu $\psi$ mit dem nach oben gerichteten
\definitionsverweis {Einheitsnormalenfeld}{}{}
und
\maabbdisp {\varphi  =
\operatorname{Id}  \times \psi} {\R^2} {Y
} {}
die Parametrisierung des Graphen. Bestimme die
\definitionsverweis {zweite Fundamentalmatrix}{}{}
zu $\varphi$.

}
{} {}

\inputaufgabe
{}
{

Wir betrachten die Einheitskugeloberfläche mit dem nach außen zeigenden
\definitionsverweis {Einheitsnormalenfeld}{}{}
und der Parametrisierung
\maabbeledisp {\varphi} {[0, 2 \pi] \times [-1,1]} {\R^3
} {(u,v)} { \left(  \sqrt{1-v^2}  \cos  u   , \,   \sqrt{1-v^2}  \sin  u  , \,   v                 \right)
} {,}
siehe
[[Kugeloberfläche/Koordinaten von Zylinder aus/Horizontale Projektion/Flächenberechnung/Beispiel|Beispiel 17.6]].
Bestimme die
\definitionsverweis {zweite Fundamentalmatrix}{}{}
zu $\varphi$.

}
{} {}

\inputaufgabe
{}
{

Wir betrachten die Einheitskugeloberfläche mit dem nach außen zeigenden
\definitionsverweis {Einheitsnormalenfeld}{}{}
und der Parametrisierung
\maabbeledisp {\varphi} {[0, 2 \pi] \times \R} {\R^3
} {(u,v)} { { \frac{   1  }{  \sqrt{1+v^2}  } }  ( \cos  u  ,  \sin  u  , v )
} {,}
siehe
[[Kugeloberfläche/Koordinaten von Zylinder aus/Mittelpunktsprojektion/Flächenberechnung/Beispiel|Beispiel 17.7]].
Bestimme die
\definitionsverweis {zweite Fundamentalmatrix}{}{}
zu $\varphi$.

}
{} {}

\inputaufgabe
{}
{

Bestimme die
\definitionsverweis {Weingartenabbildung}{}{}
bezüglich der Basis $\partial_1 \varphi, \partial_2 \varphi$ für die verschiedenen Parametrisierungen $\varphi$ der Einheitskugel aus
[[Kugeloberfläche/Koordinaten von Zylinder aus/Horizontale Projektion/Flächenberechnung/Beispiel|Beispiel 17.6]],
[[Kugeloberfläche/Koordinaten von Zylinder aus/Mittelpunktsprojektion/Flächenberechnung/Beispiel|Beispiel 17.7]]
und
[[Kugeloberfläche/Geozentrische Koordinaten/Flächenberechnung/Beispiel|Beispiel 17.8]].

}
{} {}

\inputaufgabe
{}
{

Bestimme die
\definitionsverweis {Gaußkrümmung}{}{}
der Einheitskugel mithilfe der verschiedenen Parametrisierungen $\varphi$  aus
[[Kugeloberfläche/Koordinaten von Zylinder aus/Horizontale Projektion/Flächenberechnung/Beispiel|Beispiel 17.6]],
[[Kugeloberfläche/Koordinaten von Zylinder aus/Mittelpunktsprojektion/Flächenberechnung/Beispiel|Beispiel 17.7]]
und
[[Kugeloberfläche/Geozentrische Koordinaten/Flächenberechnung/Beispiel|Beispiel 17.8]]
unter Verwendung von
[[Hyperfläche/Raum/Parametrisierung/Weingartenabbildung/Fundamentalformen/Fakt|Lemma 19.3  (4)]].

}
{} {}

\inputaufgabe
{}
{

Es sei $Y$ eine
\definitionsverweis {Kurve}{}{}
im $\R^2$. Bringe die
\definitionsverweis {Krümmung}{}{}
mit der
\definitionsverweis {zweiten Fundamentalmatrix}{}{}
in Verbindung.

}
{} {}

\inputaufgabe
{}
{

Es sei

\mavergleichskette
{\vergleichskette
{ Y
}
{ \subseteq }{ \R^3
}
{  }{
}
{    }{
}
{   }{
}
}
{}{}{}
eine
\definitionsverweis {orientierte Fläche}{}{}
und sei
\maabbdisp {\varphi} { V } { Y
} {,}

\mavergleichskette
{\vergleichskette
{ V
}
{ \subseteq }{ \R^2
}
{  }{
}
{    }{
}
{   }{
}
}
{}{}{,}
eine zweifach differenzierbare lokale Parametrisierung von $Y$ mit den Parametern $u,v$. Es sei
\mathl{\left(  g_{ij}  \right)}{} die
\definitionsverweis {erste Fundamentalmatrix}{}{}
auf $V$. Zeige, dass die
\definitionsverweis {Gaußsche Krümmung}{}{}
die Beziehung

\mavergleichskettedisp
{\vergleichskette
{ K
}
{ =} { { \frac{   1  }{  \left(  g_{11} g_{22} -g_{12}^2  \right)^2 } }  \left(  \det  \begin{pmatrix}  - { \frac{   1  }{  2 } } \partial_2^2  g_{11}  + \partial_1 \partial_2 g_{12} - { \frac{   1  }{  2 } }  \partial_1^2  g_{22}   &   { \frac{   1  }{  2 } } \partial_1 g_{11}  &  \partial_1 g_{12} - { \frac{   1  }{  2 } } \partial_2 g_{11}  \\  \partial_2 g_{12} -{ \frac{   1  }{  2 } } \partial_1 g_{22} &  g_{11}  &  g_{12}  \\ { \frac{   1  }{  2 } } \partial_2 g_{22}  &  g_{12}  &  g_{22}  \end{pmatrix} -  \det  \begin{pmatrix}  0   &   { \frac{   1  }{  2 } }  \partial_2 g_{11}  &  { \frac{   1  }{  2 } }  \partial_1 g_{22}   \\  { \frac{   1  }{  2 } } \partial_2 g_{11} & g_{11} & g_{12} \\ { \frac{   1  }{  2 } } \partial_1 g_{22}  &  g_{12} & g_{22} \end{pmatrix}  \right)
}
{ } {
}
{ } {
}
{ } {
}
}
{}{}{}
gilt.

}
{} {}

  \zwischenueberschrift{Aufgaben zum Abgeben}

\inputaufgabe
{3}
{

Es sei
\maabbeledisp {\psi} {\R^2} {\R
} {(x,y)} { x^3-2x^2y-7xy^3+x^2y^2+5y^4
} {,}
und sei $Y$ der
\definitionsverweis {Graph}{}{}
zu $\psi$ mit dem nach oben gerichteten
\definitionsverweis {Einheitsnormalenfeld}{}{}
und
\maabbdisp {\varphi  =
\operatorname{Id}  \times \psi} {\R^2} {Y
} {}
die Parametrisierung des Graphen. Bestimme die
\definitionsverweis {zweite Fundamentalmatrix}{}{}
zu $\varphi$.

}
{} {}

\inputaufgabe
{4}
{

Wir betrachten die Einheitskugeloberfläche mit dem nach außen zeigenden
\definitionsverweis {Einheitsnormalenfeld}{}{}
und der Parametrisierung
\maabbeledisp {\varphi} {{]{- { \frac{   \pi }{  2 } }},  { \frac{   \pi }{  2 } } [} \times {]{-\pi}, \pi[} } {\R^3
} {(u,v)} {(\cos  u  \cos  v  , \cos  u  \sin  v  , \sin  u )
} {,}
siehe
[[Kugeloberfläche/Geozentrische Koordinaten/Flächenberechnung/Beispiel|Beispiel 17.8]].
Bestimme die
\definitionsverweis {zweite Fundamentalmatrix}{}{}
zu $\varphi$.

}
{} {}

\inputaufgabe
{5 (2+1+1+1)}
{

Es sei
\maabb {\gamma} {]a,b[} { \R^2
} {}
eine
\definitionsverweis {differenzierbare}{}{}
\definitionsverweis {bogenparametrisierte}{}{}
injektive Kurve mit der Bildkurve

\mavergleichskettedisp
{\vergleichskette
{ C
}
{ \subseteq} { V
}
{ \subseteq} { \R^2
}
{ } {
}
{ } {
}
}
{}{}{,}
die in der offenen Teilmenge $V$ abgeschlossen sei. Wir betrachten den Zylinder über diesen Kurven, also die Abbildung
\maabbdisp {\varphi = \gamma \times
\operatorname{Id}_{ \R }} {]a,b[ \times \R } { C \times \R
} {.}
\aufzaehlungvier{Berechne die
\definitionsverweis {erste}{}{}
und die
\definitionsverweis {zweite Fundamentalmatrix}{}{}
von $\varphi$,
}{Berechne die
\definitionsverweis {Weingartenabbildung}{}{}
bezüglich der Basis $\partial_1 \varphi,  \partial_2 \varphi$.
}{Berechne die
\definitionsverweis {Hauptkrümmungen}{}{}
von $C \times \R$.
}{Berechne die
\definitionsverweis {Gaußkrümmung}{}{}
von $C \times \R$.
}

}
{} {}

 [[Kategorie:Latexseite]]
